\chapter*{Avant-propos}         % ne pas numéroter
\phantomsection\addcontentsline{toc}{chapter}{Avant-propos} % inclure dans TdM

Ce mémoire de maîtrise étant écrit, une étape notable de ma vie prend fin. Merci énormément à mon directeur de recherches Michael Lau pour avoir cru en mes capacités, pour être resté toujours positif et patient malgré cet échéancier allongé, finalement avéré nécessaire. Aussi, je lui suis très reconnaissant pour toutes ces occasions qu'il a créées ou provoquées en Amérique du Nord afin de m'aider à progresser comme apprenti mathématicien. Merci aussi au professeur Lau, pour les innombrables lettres d'appui à mes dossiers de candidatures de tous genres. Tout cela, ses contacts, ses suggestions et conseils font partie intégrante de ce qui m'a propulsé dans la situation enviable dans laquelle je me trouve en 2014-2015. J'aurai évidemment beaucoup appris à ses côté, mais j'aurai surtout gagné en assurance à l'idée de poursuivre mes études supérieures.

Un merci tout spécial au professeur Erhard Neher de l'Université d'Ottawa pour m'avoir expliqué un important détail dans une preuve de l'une de ses publications que je reprends dans ce travail. Je tiens aussi à remercier le professeur Hugo Chapdelaine de m'avoir aidé à me rendre à Paris en 2011 pour participer à un congrès international portant sur mon mathématicien favori Évariste Galois. Ce fut un séjour inspirant à plusieurs égards. Je lui suis aussi bien reconnaissant, ainsi qu'au professeur Claude Lévesque, de m'avoir occupé quelques étés de temps lors de mon passage à l'Université Laval où les quelques bourses existantes pour étudiants ne sont toujours encore que réservées aux techniciens les plus efficaces des séances d'examens. Ces marques de confiance ont été pour moi bien encourageantes. Merci à Saïd El Morchid, d'être aussi formidable, humain et pour ces fois au restaurant !

Enfin, merci à la communauté de l'anneau et aux membres du C.L.A.I.R. qui m'ont encouragé à toujours croire que j'allais bien terminer ce travail un jour. Merci aux héros de 2012 pour ces extraordinaires moments ; tous ces gens qui croient eux aussi en une éducation libre. Il est des rencontres que j'y ai faites qui m'auront assurément changé : tantôt peu, tantôt démesurément. Vive la neige et la pizza ! Je salue mon existante et unique projection ainsi que Mlle D., des personnes très généreuses que je suis content de pouvoir connaître. En terminant, je salue tout spécialement les autres mousquetaires, au nombre inhabituel de trois, avec qui j'aurai vécu non-moins d'un cinquième de mon existence avant de les quitter, eux et tous les autres, pour le Mordor. Je suis d'ailleurs bien reconnaissant à P.J. de m'y avoir accueilli si naturellement.